\begingroup
\makeatletter\tagsleft@true\let\veqno\@@leqno\makeatother
\ArithmeticHead{291}{R}{{\scriptsize ZUR THEORIE DER TRANSSCENDENTEN FUNCTIONEN. I.}}
\par\medskip\Needspace{4\baselineskip}\begin{center}[3]\end{center}
\par\indent Zieht man auf ähnliche Weise von der Reihe
\[\begin{aligned}
&\frac{1-x^{2n+2}}{1-x^{n+1}}+\frac{x^n.1-x^{2n+4}.1-x^{n+2}}{1-x^{n+1}.1-x^{n+3}}\\
&+\frac{x^{2n}.1-x^{2n+6}.1-x^{n+2}.1-x^{n+4}}{1-x^{n+1}.1-x^{n+3}.1-x^{n+5}}+\text{etc.}
\end{aligned}\]
\noindent die Reihe
\[\begin{aligned}
&\frac{x^n.1-x^{n+2}}{1-x^{n+1}}+\frac{x^{2n}.1-x^{n+2}.1-x^{n+4}}{1-x^{n+1}.1-x^{n+3}}\\
&+\frac{x^{3n}.1-x^{n+2}.1-x^{n+4}.1-x^{n+6}}{1-x^{n+1}.1-x^{n+3}.1-x^{n+5}}+\text{etc.}
\end{aligned}\]
\noindent ab, so erhält man einmal
\[\begin{aligned}
&\frac{1-x^n}{1-x^{n+1}}+\frac{x^n.1-x^n.1-x^{n+2}}{1-x^{n+1}.1-x^{n+3}}\\
&+\frac{x^{2n}.1-x^n.1-x^{n+2}.1-x^{n+4}}{1-x^{n+1}.1-x^{n+3}.1-x^{n+5}}+\text{etc.}=\psi(x,n)
\end{aligned}\]
\noindent und zweitens
\par\smallskip\hfill 37\textsuperscript{*}\par
\clearpage
\ArithmeticHead{292}{L}{ANALYSIS. NACHLASS.}
\SourceLeaf{299}
\[\begin{aligned}
&1+x^{n+1}+\frac{x^{2n+3}.1-x^{n+2}}{1-x^{n+3}}
+\frac{x^{2n+3}.x^{n+2}.1-x^{n+2}.1-x^{n+4}}{1\quad x^{n+3}.1-x^{n+5}}+\text{etc.}\\
&\hspace{2em}=1+x^{n+1}+x^{2n+3}\psi(x,n+2).
\end{aligned}\]
\noindent Man hat also, den Fall $n=0$ ausgenommen,
\[\psi(x,n)=1+x^{n+1}+x^{2n+3}\psi(x,n+2)\]
\noindent folglich
\[\psi(x,n)=1+x^{n+1}+x^{2n+3}+x^{3n+6}+x^{4n+10}+\text{etc.},\]
\noindent dagegen hat man für den Fall $n=0$
\[\psi(x,0)=1+x+x^3\psi(x,2)-\frac{1-x^2.1-x^4.1-x^6}{1-x.1-x^3.1-x^5}\ \text{etc.},\]
\noindent folglich
\begin{equation*}\tag*{[2]}
\frac{1-x^2.1-x^4.1-x^6{.}{.}{.}}{1-x.1-x^3.1-x^5{.}{.}{.}}=1+x+x^3+x^6+x^{10}+\text{etc.}
\end{equation*}
\clearpage
\endgroup
